\documentclass[10pt, letterpaper]{article}

\usepackage[top=0.4in, bottom=0.4in, left=0.6in, right=0.6in]{geometry}
\usepackage{enumitem}
\usepackage{titlesec}
\usepackage[T1]{fontenc}
\usepackage{lmodern}
\usepackage{xcolor}
\usepackage{mdframed}
\usepackage{parskip}

\definecolor{primary}{HTML}{1A3A5C}
\definecolor{accent}{HTML}{2E86AB}
\definecolor{lightbg}{HTML}{EEF4F8}
\definecolor{answerbg}{HTML}{F5F5F5}
\definecolor{starlabel}{HTML}{B94A00}

\titleformat{\section}
    {\normalsize\bfseries\color{primary}}
    {}{0em}{\MakeUppercase}
    [{\color{accent}\titlerule}]
\titlespacing{\section}{0pt}{8pt}{4pt}

\newmdenv[
    backgroundcolor=lightbg,
    linecolor=accent,
    linewidth=0.5pt,
    innerleftmargin=8pt,
    innerrightmargin=8pt,
    innertopmargin=4pt,
    innerbottommargin=4pt,
    skipabove=3pt,
    skipbelow=3pt
]{notebox}

\newmdenv[
    backgroundcolor=answerbg,
    linecolor=primary,
    linewidth=0.5pt,
    innerleftmargin=8pt,
    innerrightmargin=8pt,
    innertopmargin=6pt,
    innerbottommargin=6pt,
    skipabove=3pt,
    skipbelow=6pt
]{answerbox}

\newcommand{\q}[2]{%
    \smallskip
    \noindent\textbf{\small Q#1.}\quad\textbf{#2}
    \vspace{3pt}
}

\newcommand{\hint}[1]{%
    \begin{notebox}
    \footnotesize\textit{\textbf{Draw from:} #1}
    \end{notebox}
}

\newcommand{\s}[1]{\textbf{\textcolor{starlabel}{S:}} #1}
\newcommand{\ta}[1]{\textbf{\textcolor{starlabel}{T:}} #1}
\newcommand{\ac}[1]{\textbf{\textcolor{starlabel}{A:}} #1}
\newcommand{\res}[1]{\textbf{\textcolor{starlabel}{R:}} #1}

\pagestyle{empty}
\setlength{\parindent}{0pt}

\begin{document}

%---------- HEADER ----------
\begin{center}
    {\Large\bfseries\color{primary} Behavioral Interview Prep --- Scripted Answers}\\[2pt]
    {\small\color{accent} Jie Liu\ \ \textbar\ \ AI Product Manager\ \ \textbar\ \ JD.com / Chongbang}
\end{center}

\vspace{-2pt}
\begin{notebox}
\footnotesize\textbf{STAR Framework:} \textbf{S}ituation (context \& stakes) $\to$ \textbf{T}ask (your specific role) $\to$ \textbf{A}ction (what you did) $\to$ \textbf{R}esult (quantified impact). Aim for \textbf{2 minutes} per answer. Lead with the number, then tell the story.
\end{notebox}

%---------- LEADERSHIP ----------
\section{Leadership \& Cross-Functional Collaboration}

\q{1}{Tell me about a time you led a cross-functional team to deliver a complex project.}
\hint{JD.com: recommendation engine launch, aligned Ads, Data Analytics, Product, Engineering.}

\begin{answerbox}
\footnotesize
\s{In 2023 at JD.com, I was tasked with launching an AI-powered recommendation engine --- our biggest product bet of the year --- which required aligning engineering, data science, ads, and business teams who had never collaborated on a shared roadmap before.}

\medskip
\ta{As the AI PM, I owned the end-to-end product roadmap and was responsible for making every team clear on their role, timeline, and definition of done.}

\medskip
\ac{I kicked off with a product discovery sprint where I mapped each team's constraints and OKRs. I wrote a PRD that included not just feature requirements but model evaluation criteria and acceptance criteria in language each team could act on. I set up a shared Jira board so blockers surfaced early, and I created a single north star metric --- recommendation click-through rate --- that unified everyone's focus. When engineering and data science had a disagreement on model validation timelines, I facilitated a structured trade-off meeting and proposed a phased launch to protect both the quality bar and the commercial deadline.}

\medskip
\res{We shipped on schedule. Within one quarter the recommendation engine drove a \textbf{10\% increase in total revenue}. The process I built became the standard template for AI feature launches at JD.com.}
\end{answerbox}

\q{2}{Describe a situation where you had to resolve a conflict between teams with competing priorities.}
\hint{JD.com: engineering needed 2 extra weeks for model validation vs. business committed to a sales event go-live date.}

\begin{answerbox}
\footnotesize
\s{Mid-way through a personalization feature launch at JD.com, our engineering team flagged they needed two more weeks for model validation --- right as the business team was locked to a go-live date tied to a major sales event. Both teams had legitimate, non-negotiable priorities.}

\medskip
\ta{I was in the middle and owned the call on how to resolve it without blowing either deadline.}

\medskip
\ac{I called a structured alignment meeting and separated the technical question from the commercial one. I proposed a phased rollout: launch a rule-based fallback version on the event date to capture the commercial window, then ship the full ML model two weeks later after clean validation. I built a one-page risk matrix to show both sides the cost of each option --- delay vs. degraded model --- and walked them through it together.}

\medskip
\res{Both teams aligned in the same meeting. The fallback launched on time, protecting the sales event window. The ML version shipped two weeks later with zero quality issues. User retention metrics hit target within the first month.}
\end{answerbox}

\q{3}{Give an example of influencing stakeholders to gain buy-in on a strategy or direction.}
\hint{Chongbang: introducing AI campaign optimization tool to a resistant sales team who feared automation.}

\begin{answerbox}
\footnotesize
\s{At Chongbang, I wanted to roll out an AI-powered campaign optimization tool across all 10,000+ brand partner brands. The sales team resisted because they worried it would reduce their control and make their role feel replaceable.}

\medskip
\ta{I needed sales leadership buy-in before I could scale. I couldn't mandate it --- I had to earn it.}

\medskip
\ac{Instead of positioning the tool as a replacement, I reframed it as an amplifier of the sales team's judgment. I ran a 30-day pilot with 50 brands and built a before-and-after analysis: 20\% profitability growth in the pilot group vs. 8\% in the control group. I presented the findings to sales leadership with specific brand examples and let the data carry the argument. I also highlighted that top-performing sales reps used the tool most --- it was a differentiator within their own team.}

\medskip
\res{Leadership approved a full rollout within two weeks. Within six months the tool was live across all brand partner brands and contributed directly to our \textbf{90\% brand partner renewal rate} that year.}
\end{answerbox}

%---------- EXECUTION & RESULTS ----------
\section{Execution \& Results}

\q{4}{Walk me through a project where you drove measurable business growth.}
\hint{JD.com: 10\% revenue increase via recommendation engine and demand forecasting, driven by funnel analysis.}

\begin{answerbox}
\footnotesize
\s{In 2023 at JD.com, conversion rates had been flat for two consecutive quarters and leadership set a clear mandate: drive a step-change in revenue by end of Q3.}

\medskip
\ta{I was asked to identify and own the AI product initiatives that could move the needle.}

\medskip
\ac{I started with a deep funnel analysis in SQL and Mixpanel to find where users dropped off. The data pointed to two gaps: low relevance in search results and weak product discovery for returning users. I defined and prioritized two features --- an AI-powered semantic search upgrade and a personalized recommendation engine --- wrote the PRDs with OKRs and model evaluation criteria, and managed delivery across three teams through sprint planning. I ran weekly A/B test reviews and adjusted the algorithm ranking weights twice mid-quarter based on results.}

\medskip
\res{By end of Q3 we hit a \textbf{10\% increase in total revenue} as measured against quarterly revenue targets. The recommendation engine contributed over 60\% of that lift.}
\end{answerbox}

\q{5}{Tell me about a time you managed multiple campaigns or workstreams simultaneously.}
\hint{JD.com: personalization engine, live-streaming automation, and holiday campaign all running in parallel in Q4.}

\begin{answerbox}
\footnotesize
\s{In Q4 2023 at JD.com --- peak e-commerce season --- I was simultaneously managing three live workstreams: a personalization engine deployment, an AI live-streaming automation feature, and a major holiday campaign. All three had hard deadlines tied to the same sales window with overlapping engineering resources.}

\medskip
\ta{I had to sequence and protect all three tracks without letting any one drop, and make real-time calls when resources conflicted.}

\medskip
\ac{I built a dependency map in Jira showing which engineering and data science resources were shared across all three tracks. I negotiated a prioritization order with leadership --- personalization first as the highest revenue lever, then live-streaming, then campaign. I ran daily standups in the two weeks before launch to catch blockers same-day. When the live-streaming feature hit a data pipeline delay, I descoped one non-critical sub-feature to protect the core launch date rather than slip the whole workstream.}

\medskip
\res{All three shipped within the planned window. Combined they contributed to the \textbf{10\% revenue growth} for the quarter and the live-streaming HEART metrics came in above target.}
\end{answerbox}

\q{6}{Describe a project where you had to define scope from scratch with no prior playbook.}
\hint{Chongbang: building a brand partner growth framework for 10,000+ brands when no standard approach existed.}

\begin{answerbox}
\footnotesize
\s{When I joined Chongbang in 2020, the company had just onboarded its first wave of thousands of brand partner brands on a new global marketplace channel --- but there was no standard framework for how to grow them. I was handed a blank slate.}

\medskip
\ta{I was responsible for defining what a brand partner growth strategy looked like and building something scalable from scratch.}

\medskip
\ac{I started with voice-of-customer research --- interviewed 50+ merchants over four weeks to map their pain points and goals. I synthesized findings into three brand partner personas (high-volume/low-margin, niche/high-margin, and growth-stage) each with a tailored playbook covering ads, promotions, and event calendars. I then worked with data analytics to build KPI dashboards that tracked each brand partner's health score weekly so we could trigger interventions before churn happened rather than after.}

\medskip
\res{The framework scaled from launch to 10,000+ brands. It delivered \textbf{20\% profitability growth} for 1,000+ brands, a \textbf{90\% brand partner renewal rate}, and reduced average merchant onboarding time by \textbf{30\%}.}
\end{answerbox}

%---------- DATA & STRATEGY ----------
\section{Data-Driven Decision Making}

\q{7}{Tell me about a time you used data to change a strategy mid-execution.}
\hint{JD.com: strong top-of-funnel but declining conversion revealed wrong audience segment mid-campaign.}

\begin{answerbox}
\footnotesize
\s{Six weeks into a major user acquisition campaign at JD.com, Mixpanel data showed strong top-of-funnel numbers --- clicks and page visits were up --- but conversion to purchase was declining. Budget was already committed.}

\medskip
\ta{I had to diagnose why and make a call on whether to continue or pivot, fast.}

\medskip
\ac{I pulled a full funnel breakdown by user segment in SQL and found the campaign was bringing in high volumes of new users from a demographic with significantly lower purchase intent than our core users. I recommended narrowing the targeting parameters and reallocating 30\% of remaining budget to a re-engagement campaign targeting lapsed users --- a higher-probability conversion segment. I built a one-page data summary, presented it to the business stakeholder, and got sign-off in 24 hours.}

\medskip
\res{Within two weeks of the pivot, conversion rates recovered. We ended the campaign period with net positive results on both new user growth and conversion, and hit our quarterly KPI targets.}
\end{answerbox}

\q{8}{How have you defined and tracked KPIs for a project? Walk me through a specific example.}
\hint{JD.com: standardizing north star metrics, leading indicators, and guardrail metrics across all AI product lines.}

\begin{answerbox}
\footnotesize
\s{When I joined JD.com's AI features team, KPI definitions were inconsistent across product lines --- teams measured ``user engagement'' differently, which made joint decisions nearly impossible.}

\medskip
\ta{As PM, I took ownership of standardizing how success was defined and tracked for every feature I owned.}

\medskip
\ac{For every PRD I wrote, I defined three layers: a north star metric, two to three leading indicators, and guardrail metrics we could not let degrade. For the recommendation engine, the north star was recommendation-driven revenue, leading indicators were click-through and add-to-cart rate, and guardrails were page load time and return rate. I set up Mixpanel dashboards for each, ran weekly KPI reviews with the team, and documented every course-correction so we could learn across iterations.}

\medskip
\res{The framework was adopted across the product team. It focused our backlog prioritization on what actually moved business metrics --- which is how we concentrated effort on the initiatives that drove the \textbf{10\% revenue growth} that quarter.}
\end{answerbox}

\q{9}{Describe a time you identified a problem through data before it became a larger issue.}
\hint{Chongbang: subtle 6-week engagement decline in mid-tier brand partner cohort traced to a mis-targeted ad format.}

\begin{answerbox}
\footnotesize
\s{At Chongbang, our weekly KPI dashboard showed a cluster of mid-tier brand partners with a subtle but consistent 6-week decline in campaign engagement --- not enough to trigger an alert, but the trend line pointed toward churn, and renewal conversations were 8 weeks away.}

\medskip
\ta{I flagged the pattern and took ownership of diagnosing it before it became a retention problem.}

\medskip
\ac{I partnered with data analytics to run a deeper cohort analysis and identified the root cause: these brand partners had been auto-enrolled in a new ad format that didn't fit their product category. I worked with Customer Service to proactively reach out to the 120 affected brands with a tailored re-engagement plan --- a category-appropriate ad strategy and a one-on-one account review --- before they had any reason to consider not renewing.}

\medskip
\res{Of the 120 brands we intervened with, \textbf{105 renewed} --- an 87.5\% retention rate in a segment trending toward 60\%. This cohort was a key contributor to our overall \textbf{90\% brand partner renewal rate} that year.}
\end{answerbox}

%---------- STAKEHOLDER & RISK MANAGEMENT ----------
\section{Stakeholder \& Risk Management}

\q{10}{Tell me about a time a project faced unexpected risk or scope change. How did you handle it?}
\hint{JD.com: semantic search model failed long-tail query accuracy threshold 2 weeks before a major traffic window launch.}

\begin{answerbox}
\footnotesize
\s{Two weeks before we were scheduled to launch the AI-powered semantic search upgrade at JD.com, the data science team discovered the model's accuracy on long-tail queries --- which accounted for 35\% of our search volume --- was below the threshold defined in the PRD.}

\medskip
\ta{I had to decide: delay the launch and miss a key traffic window, or ship with a known quality gap.}

\medskip
\ac{I immediately called a risk assessment with data science, engineering, and business stakeholders. I framed three options: full delay, partial launch with a rule-based fallback for long-tail queries, or full launch with a post-launch monitoring protocol. I built a decision matrix showing the cost and risk of each option, recommended the partial launch, and communicated the revised plan to leadership the same day. I also set up an accelerated two-week sprint for data science to close the gap.}

\medskip
\res{The partial launch went live on schedule. The long-tail upgrade shipped three weeks later. Total search-driven revenue impact was within \textbf{5\% of original target} and we had zero user-facing quality complaints during the partial launch window.}
\end{answerbox}

\q{11}{Describe a situation where you had to manage a difficult stakeholder or senior decision-maker.}
\hint{Chongbang: a top-10 brand partner founder who bypassed account team and escalated directly to leadership repeatedly.}

\begin{answerbox}
\footnotesize
\s{At Chongbang, one of our largest brand partners --- a top-10 brand by revenue --- had a founder who frequently bypassed our account team and escalated directly to senior leadership whenever they disagreed with our campaign recommendations. It was damaging team morale and consuming leadership bandwidth.}

\medskip
\ta{I was brought in to manage the relationship and re-establish a structured working process without threatening the commercial relationship.}

\medskip
\ac{I requested a direct meeting with the founder --- not to defend past decisions but to genuinely understand their frustrations. I listened first. Then I walked through our performance data for their brand specifically and showed the decisions that had driven their results. I proposed a monthly business review cadence with direct KPI visibility and a clear channel for raising concerns before escalating. I also assigned them a dedicated point of contact on my team so there was always someone accountable.}

\medskip
\res{Escalations stopped within six weeks. The brand \textbf{renewed at a higher contract tier}. The monthly review format we built for them was later rolled out to our top 50 brand partners as a standard practice.}
\end{answerbox}

\q{12}{Give an example of when you had to say no --- or push back --- to protect a project's scope or quality.}
\hint{JD.com: marketing requested out-of-scope personalization feature two weeks before a live-streaming launch.}

\begin{answerbox}
\footnotesize
\s{Two weeks before a major live-streaming feature launch at JD.com, the marketing team requested an additional personalization layer --- a genuinely good idea, but out of scope and requiring two more weeks of engineering work.}

\medskip
\ta{I had to push back without damaging the relationship and without slipping the launch date.}

\medskip
\ac{I set up a quick call with the marketing lead, acknowledged the value of the idea directly, and then showed a one-page impact analysis: a two-week delay would cost us the peak traffic window, which translated to a meaningful estimated revenue loss. I proposed we formally log it as a high-priority item in the next sprint backlog and commit to a specific delivery date --- four weeks post-launch. I documented the agreement in writing and followed up with a Jira ticket so it didn't get lost.}

\medskip
\res{The original launch shipped on time. The personalization layer shipped four weeks later as promised. The marketing lead later told me it was the first time a PM had pushed back on a request and then actually followed through on the committed timeline.}
\end{answerbox}

%---------- GROWTH & ADAPTABILITY ----------
\section{Growth \& Adaptability}

\q{13}{Tell me about a time you had to learn something new quickly to get a project done.}
\hint{JD.com: joining an ML team without a data science background, needing to earn credibility to write effective PRDs.}

\begin{answerbox}
\footnotesize
\s{When I joined JD.com as an AI PM, I had strong project management and campaign operations experience but had never worked directly with data science or ML teams. I was expected to write PRDs that included model evaluation criteria and speak credibly to engineers about feature design.}

\medskip
\ta{I needed to build enough ML fluency to earn the team's trust and communicate effectively --- not to become a data scientist, but to stop being a bottleneck.}

\medskip
\ac{I completed Andrew Ng's Machine Learning Specialization on Coursera in six weeks while working full-time, focusing specifically on the concepts that came up in our sprint reviews: precision/recall tradeoffs, overfitting, and recommendation system architectures. I also paired with our data scientists during model review sessions, asking questions about how their outputs mapped to user experience. Within a month, the team told me my PRDs were the clearest technical-to-product translations they had worked from.}

\medskip
\res{Fewer revision cycles on PRDs meant \textbf{faster sprint velocity} across the team. It also directly shaped how I scoped the recommendation engine --- which contributed to the \textbf{10\% revenue growth} we hit that quarter.}
\end{answerbox}

\q{14}{Describe a project that didn't go as planned. What would you do differently?}
\hint{Chongbang: peak event campaign underperformed due to an unmonitored platform algorithm change.}

\begin{answerbox}
\footnotesize
\s{At Chongbang, I led a promotional campaign for 500 brands during a peak sales event. I built the strategy on prior performance data and set aggressive targets --- but didn't account for a platform algorithm change that reduced organic reach the week before the event.}

\medskip
\ta{The campaign underperformed. Several brands missed their sales targets and I owned the outcome.}

\medskip
\ac{I ran a post-mortem with the data and account teams. The algorithm change had been announced internally two weeks prior --- I hadn't built a process to monitor platform changes that could affect campaign inputs. I had also relied too heavily on historical benchmarks without building scenario planning for external variables. I rebuilt our campaign planning process to include a pre-launch risk checklist and a weekly platform change monitor.}

\medskip
\res{That campaign missed target by roughly \textbf{15\%}. The following event cycle, with the new process in place, we \textbf{exceeded target by 12\%}. The lesson I carry: good historical data is necessary but not sufficient --- you also need a process for identifying what the data can't see.}
\end{answerbox}

\q{15}{Where do you see yourself in 3--5 years, and how does this role fit that path?}
\hint{Senior AI PM owning model-to-business strategy, building on current MS program and JD.com AI product experience.}

\begin{answerbox}
\footnotesize
My goal is to grow into a senior AI PM role at a company where machine learning is core to the user experience --- not bolted on. In the near term, I'm deepening my foundation through my MS in Project Management at Northeastern, where I'm focused on risk management and analytics-driven leadership.

\medskip
At the same time, I'm actively building my understanding of LLM product development --- how to define evaluation frameworks for generative AI features, how to manage the model-to-product feedback loop, and how to structure roadmaps around AI capabilities that are still maturing.

\medskip
In 3--5 years, I see myself owning a product area where I'm defining not just the feature roadmap but the AI strategy --- what problems we solve with ML, how we measure model-to-business impact, and how we build trust with users around AI-powered experiences.

\medskip
This role fits that path because the product challenges here require exactly the skills I've been building: cross-functional execution on ML features, data-driven roadmap prioritization, and the ability to translate between technical and business teams at speed.
\end{answerbox}

%---------- QUESTIONS TO ASK ----------
\section{Questions to Ask the Interviewer}

\begin{enumerate}[leftmargin=1.5em, itemsep=3pt]
    \item What does success look like for this role in the first 90 days?
    \item How does the PM team work with data science and engineering day-to-day --- are PMs expected to write specs at the model level?
    \item What are the biggest AI product challenges the team is currently navigating?
    \item How does the organization decide which problems to solve with AI vs. traditional approaches?
    \item What does the growth path look like for AI PMs here?
\end{enumerate}

\end{document}
